\documentclass[12pt,a4paper]{article}
\usepackage[T2A]{fontenc}
\usepackage[utf8]{inputenc}
\usepackage[russian]{babel}
\usepackage{float}
\usepackage{amsmath, amssymb}
\usepackage{geometry}
\usepackage{tikz}
\usepackage{tikz-3dplot}
\usepackage{caption}
\usetikzlibrary{arrows.meta, calc, decorations.pathmorphing}
\geometry{margin=2cm}

\newcommand{\vu}{\bar u}
\newcommand{\vv}{\bar v}
\newcommand{\vx}{\bar x}
\newcommand{\vy}{\bar y}
\newcommand{\ve}{\bar e}
\newcommand{\vo}{\bar o}

\begin{document}

\setcounter{page}{112}

\begin{center}
  \textbf{ЛЕКЦИЯ 16. РАЗМЕРНОСТЬ И БАЗИС ЛИНЕЙНОГО ПРОСТРАНСТВА}
\end{center}

\subsection*{\textbf{16.1. Понятие базиса и размерности}}

Говорят, что линейное пространство $V$ имеет размерность $n$, если в
пространстве $V$ существует $n$ линейно независимых векторов, а любая
система из $n + 1$ вектора линейно зависима.

\textbf{Определение.} Размерность линейного пространства --- это
наибольшее число линейно независимых векторов в этом пространстве.

\textbf{Обозначение.} $\dim V$ --- размерность линейного пространства $V$.

\textit{Первое определение базиса}

\textbf{Определение.} Пусть $V$ --- $n$-мерное линейное пространство.

Базисом $\mathcal{B} = \{\ve_1, \ve_2, \ldots, \ve_n\} \subset V$
называется любая система из $n$ линейно независимых векторов в этом
пространстве.

\textbf{Пример}

В линейном пространстве матриц размера $2 \times 2$ ---
$\mathbb{R}^{2\times 2}$ рассмотрим векторы (матрицы):
$$
\ve_1 =
\begin{pmatrix} 1 & 0 \\ 0 & 0
\end{pmatrix}; \quad \ve_2 =
\begin{pmatrix} 0 & 1 \\ 0 & 0
\end{pmatrix}; \quad \ve_3 =
\begin{pmatrix} 0 & 0 \\ 1 & 0
\end{pmatrix}; \quad \ve_4 =
\begin{pmatrix} 0 & 0 \\ 0 & 1
\end{pmatrix}.
$$

Векторы $\ve_1, \ve_2, \ve_3, \ve_4$ --- линейно независимы, т.к.
$$
\lambda_1 \ve_1 + \lambda_2 \ve_2 + \lambda_3 \ve_3 + \lambda_4 \ve_4 =
\begin{pmatrix} \lambda_1 & \lambda_2 \\ \lambda_3 & \lambda_4
\end{pmatrix} = O =
\begin{pmatrix} 0 & 0 \\ 0 & 0
\end{pmatrix} \Leftrightarrow
$$
$$
\Leftrightarrow \lambda_1 = \lambda_2 = \lambda_3 = \lambda_4 = 0.
$$

Любые 5 матриц размера $2 \times 2$ линейно зависимы

$\Rightarrow \dim \mathbb{R}^{2\times 2} = 4$ и $\mathcal{B} =
\{\ve_1, \ve_2, \ve_3, \ve_4\}$ --- базис

\textbf{Пример}

В линейном пространстве многочленов степени $\le 2$ ---

$P_2 = \{p(t) = at^2 + bt + c; a, b, c \in \mathbb{R}\}$ рассмотрим многочлены:

$t^2, t, 1$ --- линейно независимы, т.к.
$$
\lambda_1 t^2 + \lambda_2 t + \lambda_3 \cdot 1 = 0 \Leftrightarrow
\lambda_1 = \lambda_2 = \lambda_3 = 0.
$$

Любые 4 многочлена линейно зависимы

$\Rightarrow \dim P_2 = 3$ и $\mathcal{B} = \{t^2, t, 1\}$ --- базис

\subsection*{\textbf{16.2. Теорема о разложении вектора по базису.
Преобразование координат вектора при переходе к другому базису}}

\textbf{Теорема.} (О разложении вектора по базису).

Пусть $V$ --- линейное пространство и $\mathcal{B} = \{\ve_1, \ve_2,
\ldots, \ve_n\}$ --- базис в $V$. Тогда любой вектор $\vu \in V$ есть
линейная комбинация векторов базиса, т.е.

$\forall \vu \in V$ $\exists$ числа $x_1, x_2, \ldots, x_n$ такие, что
\begin{equation}
  \vu = x_1 \ve_1 + x_2 \ve_2 + \cdots + x_n \ve_n \tag{16.1}
\end{equation}
и разложение (16.1) единственно.

\textbf{Доказательство}

1) Докажем существование разложения (16.1)

Возьмем любой вектор $\vu \in V$. Рассмотрим систему векторов $\Sigma
= \{\vu, \ve_1, \ve_2, \ldots, \ve_n\}$. Система $\Sigma$ состоит из
$n + 1$ вектора. Т.к. $\dim V = n$, то $\Sigma$ --- линейно зависима,
т.е. существует нетривиальная линейная комбинация
\begin{equation}
\lambda_0 \vu + \lambda_1 \ve_1 + \lambda_2 \ve_2 + \cdots +
\lambda_n \ve_n = \vo, \tag{16.2}
\end{equation}
где хотя бы один из коэффициентов $\lambda_0, \lambda_1, \lambda_2,
\ldots, \lambda_n$ отличен от нуля.

Предположим, что $\lambda_0 = 0 \Rightarrow \lambda_1 \ve_1 +
\lambda_2 \ve_2 + \cdots + \lambda_n \ve_n = \vo$

Но $\ve_1, \ve_2, \ldots, \ve_n$ --- линейно независимы

$\Rightarrow \lambda_1 = \lambda_2 = \cdots = \lambda_n = 0
\Rightarrow$ противоречие и $\lambda_0 \ne 0$

Из равенства (16.2) получим:
$$
\vu = -\frac{\lambda_1}{\lambda_0} \ve_1 -
\frac{\lambda_2}{\lambda_0} \ve_2 - \cdots - \frac{\lambda_n}{\lambda_0} \ve_n =
$$
$$
= x_1 \ve_1 + x_2 \ve_2 + \cdots + x_n \ve_n
$$
--- разложение (16.1).

2) Докажем единственность разложения (16.1)

Пусть существуют два различных разложения:
$$
\vu = x_1 \ve_1 + x_2 \ve_2 + \cdots + x_n \ve_n
$$
и
\begin{equation}
\vu = y_1 \ve_1 + y_2 \ve_2 + \cdots + y_n \ve_n. \tag{16.3}
\end{equation}

Вычитаем из равенства (16.1) равенство (16.3):
\begin{equation}
\vo = (x_1 - y_1) \ve_1 + (x_2 - y_2) \ve_2 + \cdots + (x_n - y_n)
\ve_n. \tag{16.4}
\end{equation}

Т.к. $\ve_1, \ve_2, \ldots, \ve_n$ --- линейно независимы, то
линейная комбинация (16.4) тривиальна

$\Rightarrow x_1 - y_1 = 0, x_2 - y_2 = 0, \ldots, x_n - y_n = 0$

$\Rightarrow x_1 = y_1, x_2 = y_2, \ldots, x_n = y_n \Rightarrow$ противоречие

$\Rightarrow$ разложения (16.1) единственно.

\textbf{ч.т.д.}

Замечание. Числа $x_1, x_2, \ldots, x_n$ называются координатами
вектора $\vu$ в базисе $\mathcal{B}$.

\textbf{Обозначение.} Если $\mathcal{B} = \{\ve_1, \ve_2, \ldots,
\ve_n\}$ --- базис в $V$, то
$$
\vu =
\begin{pmatrix} x_1 \\ x_2 \\ \vdots \\ x_n
\end{pmatrix}_{\mathcal{B}} \Leftrightarrow \vu = x_1 \ve_1 + x_2
\ve_2 + \cdots + x_n \ve_n.
$$

Пусть $V$ --- линейное пространство и $\mathcal{B} = \{\ve_1, \ve_2,
\ldots, \ve_n\}$ --- базис в $V$.

\textbf{Теорема.} Если $\vu =
\begin{pmatrix} x_1 \\ x_2 \\ \vdots \\ x_n
\end{pmatrix}_{\mathcal{B}}$ и $\vv =
\begin{pmatrix} y_1 \\ y_2 \\ \vdots \\ y_n
\end{pmatrix}_{\mathcal{B}}$, то
$$
\vu + \vv =
\begin{pmatrix} x_1 + y_1 \\ x_2 + y_2 \\ \vdots \\ x_n + y_n
\end{pmatrix}_{\mathcal{B}} \quad \text{и} \quad \lambda \vu =
\begin{pmatrix} \lambda x_1 \\ \lambda x_2 \\ \vdots \\ \lambda x_n
\end{pmatrix}_{\mathcal{B}}.
$$

\textbf{Доказательство}
$$
\vu = \sum_{i=1}^n x_i \ve_i; \quad \vv = \sum_{i=1}^n y_i \ve_i
\Rightarrow \vu + \vv = \sum_{i=1}^n x_i \ve_i + \sum_{i=1}^n y_i
\ve_i = \sum_{i=1}^n (x_i + y_i) \ve_i
$$
$$
\Rightarrow \vu + \vv =
\begin{pmatrix} x_1 + y_1 \\ x_2 + y_2 \\ \vdots \\ x_n + y_n
\end{pmatrix}_{\mathcal{B}}.
$$

Аналогично для $\lambda \vu$:
$$
\lambda \vu = \lambda \sum_{i=1}^n x_i \ve_i = \sum_{i=1}^n \lambda
x_i \ve_i \Rightarrow \lambda \vu =
\begin{pmatrix} \lambda x_1 \\ \lambda x_2 \\ \vdots \\ \lambda x_n
\end{pmatrix}_{\mathcal{B}}.
$$

\textbf{ч.т.д.}

\textit{Преобразование координат вектора при переходе к другому базису}

Пусть $V$ --- линейное пространство и $\mathcal{B} = \{\ve_1, \ve_2,
\ldots, \ve_n\}$ и $\mathcal{B}' = \{\ve'_1, \ve'_2, \ldots,
\ve'_n\}$ --- базисы в $V$

$\mathcal{B}$ --- исходный базис; $\mathcal{B}'$ --- новый базис

Пусть $\vu =
\begin{pmatrix} x_1 \\ x_2 \\ \vdots \\ x_n
\end{pmatrix}_{\mathcal{B}}$ и $\vu =
\begin{pmatrix} x'_1 \\ x'_2 \\ \vdots \\ x'_n
\end{pmatrix}_{\mathcal{B}'}$

Пусть векторы нового базиса раскладываются по векторам старого базиса:
\begin{equation}
\begin{cases}
\ve'_1 = c_{11} \ve_1 + c_{21} \ve_2 + \cdots + c_{n1} \ve_n;\\
\ve'_2 = c_{12} \ve_1 + c_{22} \ve_2 + \cdots + c_{n2} \ve_n;\\
\cdots\\
\ve'_i = c_{1i} \ve_1 + c_{2i} \ve_2 + \cdots + c_{ni} \ve_n;\\
\cdots\\
\ve'_n = c_{1n} \ve_1 + c_{2n} \ve_2 + \cdots + c_{nn} \ve_n,
\end{cases}
\tag{16.5}
\end{equation}
$$
\Rightarrow \ve'_i =
\begin{pmatrix} c_{1i} \\ c_{2i} \\ \vdots \\ c_{ni}
\end{pmatrix}_{\mathcal{B}}.
$$

Введем матрицу
$$
C =
\begin{pmatrix} c_{11} & c_{12} & \ldots & c_{1n} \\ c_{21} & c_{22}
& \ldots & c_{2n} \\ \ldots & \ldots & \ldots & \ldots \\ c_{n1} &
c_{n2} & \ldots & c_{nn}
\end{pmatrix}
$$
--- матрица перехода от базиса $\mathcal{B}$ к базису $\mathcal{B}'$.

Равенство (16.5) можно записать в матричном виде:
\begin{equation}
\mathcal{E}' = \mathcal{E} \cdot C, \tag{16.6}
\end{equation}
где $\mathcal{E} = (\ve_1, \ve_2, \ldots, \ve_n)$; $\mathcal{E}' =
(\ve'_1, \ve'_2, \ldots, \ve'_n)$.

Т.к. $\vu =
\begin{pmatrix} x'_1 \\ x'_2 \\ \vdots \\ x'_n
\end{pmatrix}_{\mathcal{B}'}$; $\ve'_i =
\begin{pmatrix} c_{1i} \\ c_{2i} \\ \vdots \\ c_{ni}
\end{pmatrix}_{\mathcal{B}}$ и $\vu =
\begin{pmatrix} x_1 \\ x_2 \\ \vdots \\ x_n
\end{pmatrix}_{\mathcal{B}}$
$$
\Rightarrow \vu = \sum_{i=1}^n x'_i \ve'_i = \sum_{i=1}^n x'_i
\sum_{j=1}^n c_{ji} \ve_j = \sum_{j=1}^n \left( \sum_{i=1}^n x'_i
c_{ji} \right) \ve_j.
$$

С другой стороны
$$
\vu = \sum_{j=1}^n x_j \ve_j.
$$

В силу единственности разложения вектора по базису получим
\begin{equation}
x_j = \sum_{i=1}^n x'_i c_{ji}, \quad j = 1, 2, \ldots, n \tag{16.7}
\end{equation}

Запишем равенства (16.7) подробно

\begin{equation}
\begin{cases}
x_1 = c_{11} x'_1 + c_{12} x'_2 + \cdots + c_{1n} x'_n;\\
x_2 = c_{21} x'_1 + c_{22} x'_2 + \cdots + c_{2n} x'_n;\\
\cdots\\
x_j = c_{j1} x'_1 + c_{j2} x'_2 + \cdots + c_{jn} x'_n;\\
\cdots\\
x_n = c_{n1} x'_1 + c_{n2} x'_2 + \cdots + c_{nn} x'_n.
\end{cases}
\tag{16.8}
\end{equation}

Запишем систему (16.8) в матричном виде:
\begin{equation}
\begin{pmatrix} x_1 \\ x_2 \\ \vdots \\ x_n
\end{pmatrix}_{\mathcal{B}} = C \cdot
\begin{pmatrix} x'_1 \\ x'_2 \\ \vdots \\ x'_n
\end{pmatrix}_{\mathcal{B}'}, \tag{16.9}
\end{equation}
где
$$
C =
\begin{pmatrix} c_{11} & c_{12} & \ldots & c_{1n} \\ c_{21} & c_{22}
& \ldots & c_{2n} \\ \ldots & \ldots & \ldots & \ldots \\ c_{n1} &
c_{n2} & \ldots & c_{nn}
\end{pmatrix}
$$
--- матрица перехода от базиса $\mathcal{B}$ к базису $\mathcal{B}'$.

\textbf{Теорема.} Матрица перехода от базиса $\mathcal{B}$ к базису
$\mathcal{B}'$ не вырождена.

\textbf{Доказательство}

Используем формулу (16.6):
$$
\mathcal{E}' = \mathcal{E} \cdot C,
$$
где $\mathcal{E} = (\ve_1, \ve_2, \ldots, \ve_n)$; $\mathcal{E}' =
(\ve'_1, \ve'_2, \ldots, \ve'_n)$
$$
\Rightarrow \mathcal{E} = \mathcal{E}' \cdot C' = \mathcal{E} \cdot C
\cdot C' = \mathcal{E} \cdot D, \text{ где}
$$
$$
D = C \cdot C'.
$$

Покажем, что $D = (d_{ij}) = E$ --- единичная матрица
$$
\ve_i = \ve_1 \cdot d_{1i} + \ve_2 \cdot d_{2i} + \cdots + \ve_i
\cdot d_{ii} + \cdots + \ve_n \cdot d_{ni}
$$
\begin{equation}
\Rightarrow \vo = \ve_1 \cdot d_{1i} + \cdots + \ve_i \cdot (d_{ii} -
1) + \cdots + \ve_n \cdot d_{ni}. \tag{16.10}
\end{equation}

Т.к. векторы $\ve_1, \ldots, \ve_n$ --- линейно независимы

$\Rightarrow$ линейная комбинация (16.10) --- тривиальная
$$
\Rightarrow d_{1i} = 0, \ldots, d_{ii} - 1 = 0, \ldots, d_{ni} = 0
$$
$$
\Rightarrow d_{ij} = \delta_{ij} =
\begin{cases} 1, & \text{при } i = j \\ 0, & \text{при } i \ne j
\end{cases}
$$
--- символ Кронекера
$$
\Rightarrow D = C \cdot C' = E \Rightarrow C' = C^{-1}
$$
$\Rightarrow C$ --- не вырождена.

\textbf{ч.т.д.}

Т.к. $C$ --- не вырождена, то из формулы (16.9) получим:

\begin{equation}
\begin{pmatrix} x'_1 \\ x'_2 \\ \vdots \\ x'_n
\end{pmatrix}_{\mathcal{B}'} = C^{-1} \cdot
\begin{pmatrix} x_1 \\ x_2 \\ \vdots \\ x_n
\end{pmatrix}_{\mathcal{B}} \tag{16.11}
\end{equation}

(16.11) --- преобразование координат вектора при переходе к новому базису.

\subsection*{\textbf{16.3. Линейная оболочка. Система образующих
линейного пространства. Второе определение базиса}}

\textit{Линейная оболочка}

Пусть $\Sigma = \{\vu_1, \vu_2, \ldots, \vu_m\}$ --- система векторов
линейного пространства $V$.

\textbf{Определение.} Линейной оболочкой системы векторов $\Sigma =
\{\vu_1, \vu_2, \ldots, \vu_m\}$ называется множество всевозможных
линейных комбинаций этих векторов.

\textbf{Обозначение.} $L(\vu_1, \vu_2, \ldots, \vu_m)$ --- линейная
оболочка системы векторов $\Sigma = \{\vu_1, \vu_2, \ldots, \vu_m\}$.

Таким образом,
$$
L(\vu_1, \vu_2, \ldots, \vu_m) = \{\vu = \lambda_1 \vu_1 + \lambda_2
\vu_2 + \cdots + \lambda_m \vu_m; \lambda_i \in \mathbb{R}; i = 1, \ldots, m\}
$$

\textbf{Теорема.} Если $\vu_1, \vu_2, \ldots, \vu_m$ --- векторы
линейного пространства $V$, то $L(\vu_1, \vu_2, \ldots, \vu_m)$
является линейным подпространством пространства $V$.

Теорема вытекает из определения линейного подпространства.

\textbf{Теорема.} Размерность линейного подпространства не
превосходит размерности пространства. Линейное подпространство той же
размерности, что и все пространство, совпадает с пространством.

\textbf{Доказательство.}

Первая часть формулировки теоремы очевидна.

Пусть $M$ --- линейное подпространство в линейном пространстве $V$
$\Rightarrow M \subset V$.

Пусть $\dim M = \dim V = n$.

Возьмем в $M$ линейно независимую систему векторов $\Sigma = \{\ve_1,
\ve_2, \ldots, \ve_n\}$.

Система векторов $\Sigma$ является базисом для обоих пространств $M$ и $V$.
$$
\Rightarrow \forall \vu \in V \Rightarrow \vu = x_1 \ve_1 + x_2 \ve_2
+ \cdots + x_n \ve_n \Rightarrow \vu \in M
$$
$$
\Rightarrow V \subset M \Rightarrow M = V
$$

\textbf{ч.т.д.}

\textbf{Теорема.} Если $\vx = \lambda_1 \vu_1 + \lambda_2 \vu_2 +
\cdots + \lambda_m \vu_m$, то
$$
L(\vx, \vu_1, \vu_2, \ldots, \vu_m) = L(\vu_1, \vu_2, \ldots, \vu_m).
$$

\textbf{Доказательство}

Очевидно, что $L(\vu_1, \vu_2, \ldots, \vu_m) \subset L(\vx, \vu_1,
\vu_2, \ldots, \vu_m)$

$\forall \vv \in L(\vx, \vu_1, \vu_2, \ldots, \vu_m) \Rightarrow$
$$
\Rightarrow \vv = \alpha_0 \vx + \alpha_1 \vu_1 + \alpha_2 \vu_2 +
\cdots + \alpha_m \vu_m =
$$
$$
= \alpha_0 (\lambda_1 \vu_1 + \lambda_2 \vu_2 + \cdots + \lambda_m \vu_m) +
$$
$$
+ \alpha_1 \vu_1 + \alpha_2 \vu_2 + \cdots + \alpha_m \vu_m =
$$
$$
= \beta_1 \vu_1 + \beta_2 \vu_2 + \cdots + \beta_m \vu_m \in L(\vu_1,
\vu_2, \ldots, \vu_m)
$$
$$
\Rightarrow L(\vx, \vu_1, \vu_2, \ldots, \vu_m) \subset L(\vu_1,
\vu_2, \ldots, \vu_m)
$$
$$
\Rightarrow L(\vx, \vu_1, \vu_2, \ldots, \vu_m) = L(\vu_1, \vu_2,
\ldots, \vu_m).
$$

\textbf{ч.т.д.}

\textbf{Определение.} Cистема $\Sigma = \{\vu_1, \vu_2, \ldots,
\vu_m\}$ называется системой образующих линейного пространства $V$,
если ее линейная оболочка совпадает с пространством: $L(\vu_1, \vu_2,
\ldots, \vu_m) = V$.

\textbf{Теорема.} (Теорема Штейница).

Пусть $\Sigma = \{\vu_1, \vu_2, \ldots, \vu_m\}$ --- система
образующих линейного пространства $V$ и $\Sigma_1 = \{\vv_1, \vv_2,
\ldots, \vv_n\}$ --- линейно независимая система в пространстве $V$.

Тогда

1) $m \ge n$;

2) какие-то $n$ векторов системы $\Sigma$ можно заменить на векторы
системы $\Sigma_1$ так, что полученная система останется системой
образующих пространства $V$.

\textbf{Теорема.} Любая линейно независимая система образующих
линейного пространства $V$ является его базисом.

\textbf{Доказательство}

Пусть $\mathcal{B} = \{\ve_1, \ve_2, \ldots, \ve_n\}$ --- линейно
независимая система образующих пространства $V$. Покажем, что $\dim V = n$.

В пространстве $V$ существует $n$ линейно независимых векторов.
Покажем, что любые $n+1$ векторов линейно зависимы.

Берем любую систему из $n+1$ вектора $\Sigma = \{\vu_1, \vu_2,
\ldots, \vu_{n+1}\}$. Предположим, что $\Sigma$ --- линейно независима.

Но т.к. $\mathcal{B} = \{\ve_1, \ve_2, \ldots, \ve_n\}$ --- система
образующих пространства $V \Rightarrow$ по теореме Штейница $n \ge n + 1$

$\Rightarrow$ противоречие $\Rightarrow \Sigma$ --- линейно зависима

$\Rightarrow \dim V = n$ и $\mathcal{B} = \{\ve_1, \ve_2, \ldots,
\ve_n\}$ --- базис в $V$

\textbf{ч.т.д.}

\textit{Второе определение базиса}

\textbf{Определение.} $\mathcal{B} = \{\ve_1, \ve_2, \ldots, \ve_n\}$
--- базис в линейном пространстве $V$, если

1) $\mathcal{B} = \{\ve_1, \ve_2, \ldots, \ve_n\}$ --- линейно независима;

2) $L(\mathcal{B}) = V$, т.е. $\mathcal{B} = \{\ve_1, \ve_2, \ldots,
\ve_n\}$ --- система образующих пространства $V$.

\textit{Определение базиса, удобное для решения задач}

\textbf{Определение.} $\mathcal{B} = \{\ve_1, \ve_2, \ldots, \ve_n\}$
--- базис в линейном пространстве $V$, если

1) $\ve_1, \ve_2, \ldots, \ve_n$ --- линейно независимы;

2) $\forall \vx \in V \Rightarrow \vx = x_1 \ve_1 + x_2 \ve_2 +
\cdots + x_n \ve_n$.

\textbf{Пример.} Пусть $M$ --- множество многочленов степени $\le 4$,
которые делятся на $t^2 + t + 1$. Известно, что $M$ --- линейное
подпространство в линейном пространстве $P_4$. Найти базис и
размерность пространства $M$.

\textbf{Решение}
$$
M = \{p(t) = (at^2 + bt + c)(t^2 + t + 1); a, b, c \in \mathbb{R}\}
$$
$$
p(t) = (at^2 + bt + c)(t^2 + t + 1) =
$$
$$
= a(t^4 + t^3 + t^2) + b(t^3 + t^2 + t) + c(t^2 + t + 1)
$$

1) Рассмотрим векторы (многочлены):
$$
\ve_1 = (t^4 + t^3 + t^2); \quad \ve_2 = t^3 + t^2 + t; \quad \ve_3 =
t^2 + t + 1
$$
$$
\lambda_1 \ve_1 + \lambda_2 \ve_2 + \lambda_3 \ve_3 = 0
\Leftrightarrow \lambda_1 = \lambda_2 = \lambda_3 = 0
$$
$\Rightarrow \ve_1, \ve_2, \ve_3$ --- линейно независимы

2) $\forall p(t) \in M \Rightarrow p(t) = a \ve_1 + b \ve_2 + c \ve_3$
$$
\Rightarrow \mathcal{B} = \{\ve_1, \ve_2, \ve_3\} =
$$
$$
= \{t^4 + t^3 + t^2; t^3 + t^2 + t; t^2 + t + 1\}
$$
--- базис в $M$
$$
\Rightarrow \dim M = 3
$$

\textbf{Пример.} Пусть $M$ --- множество симметричных матриц размера
$2 \times 2$. Известно, что $M$ --- линейное подпространство в
линейном пространстве $\mathbb{R}^{2\times 2}$. Найти базис и
размерность пространства $M$.

\textbf{Решение}
$$
M = \left\{X =
\begin{pmatrix} a & c \\ c & b
\end{pmatrix}; a, b, c \in \mathbb{R}\right\};
$$
$$
X =
\begin{pmatrix} a & c \\ c & b
\end{pmatrix} = a
\begin{pmatrix} 1 & 0 \\ 0 & 0
\end{pmatrix} + b
\begin{pmatrix} 0 & 0 \\ 0 & 1
\end{pmatrix} + c
\begin{pmatrix} 0 & 1 \\ 1 & 0
\end{pmatrix}.
$$

1) Рассмотрим векторы (матрицы):
$$
\ve_1 =
\begin{pmatrix} 1 & 0 \\ 0 & 0
\end{pmatrix}; \quad \ve_2 =
\begin{pmatrix} 0 & 0 \\ 0 & 1
\end{pmatrix}; \quad \ve_3 =
\begin{pmatrix} 0 & 1 \\ 1 & 0
\end{pmatrix};
$$
$$
\lambda_1 \ve_1 + \lambda_2 \ve_2 + \lambda_3 \ve_3 =
$$
$$
= \lambda_1
\begin{pmatrix} 1 & 0 \\ 0 & 0
\end{pmatrix} + \lambda_2
\begin{pmatrix} 0 & 0 \\ 0 & 1
\end{pmatrix} + \lambda_3
\begin{pmatrix} 0 & 1 \\ 1 & 0
\end{pmatrix} =
$$
$$
=
\begin{pmatrix} \lambda_1 & \lambda_3 \\ \lambda_3 & \lambda_2
\end{pmatrix} = O =
\begin{pmatrix} 0 & 0 \\ 0 & 0
\end{pmatrix} \Leftrightarrow \lambda_1 = \lambda_2 = \lambda_3 = 0,
$$

$\Rightarrow \ve_1, \ve_2, \ve_3$ --- линейно независимы.

2) $\forall X \in M \Rightarrow X = a \ve_1 + b \ve_2 + c \ve_3$
$$
\Rightarrow \mathcal{B} = \{\ve_1, \ve_2, \ve_3\} =
$$
$$
= \left\{
\begin{pmatrix} 1 & 0 \\ 0 & 0
\end{pmatrix},
\begin{pmatrix} 0 & 0 \\ 0 & 1
\end{pmatrix},
\begin{pmatrix} 0 & 1 \\ 1 & 0
\end{pmatrix} \right\}
$$
--- базис в $M$
$$
\Rightarrow \dim M = 3.
$$

\newpage

\begin{center}
\textbf{СПИСОК ЛИТЕРАТУРЫ}
\end{center}

1. Лившиц К.И. Курс линейной алгебры и аналитической геометрии
[Электронный ресурс]: Санкт-Петербург: Лань, 2017.-508 с.---Режим
доступа: https://e.lanbook.com/book/93697.

2. Горлач Б.А. Линейная алгебра и аналитическая геометрия
[Электронный ресурс]: Санкт-Петербург: Лань, 2017.-300 с.---Режим
доступа: https://e.lanbook.com/book/99103.

3. Клетеник Д.В. Сборник задач по аналитической геометрии: учебное
пособие для втузов.---СПб: Профессия, 2006.---199 с.

4. Проскуряков И.В. Сборник задач по линейной алгебре [Электронный
ресурс]: учебное пособие.---Санкт-Петербург: Лань, 2019.-476
с.---Режим доступа: https://e.lanbook.com/book/114701.

5. Бутузов В.Ф., Крутицкая Н.Ч., Шишкин А.А. Линейная алгебра в
вопросах и задачах: Рек. Минобрнауки РФ в качестве учебного пособия
для втузов.---М.:Лань, 2006.---247 с.

6. Беклемишев Д.В. Курс аналитической геометрии и линейной алгебры
[Электронный ресурс]: учебник.---Санкт-Петербург: Лань, 2020.-448
с.---Режим доступа: https://e.lanbook.com/book/126146.

7. Курош А.Г. Курс высшей алгебры [Электронный ресурс]:
учебник.---Санкт-Петербург: Лань, 2019.-476 с.---Режим доступа:
https://e.lanbook.com/book/114701.

8. Беклемишева Л.А., Беклемишев Д.В., Петрович А.Ю., Чубаров И.А.
Сборник задач по аналитической геометрии и линейной алгебре
[Электронный ресурс]: учебное пособие.---Санкт-Петербург: Лань,
2019.-496 с.---Режим доступа: https://e.lanbook.com/book/122183.

\end{document}
